\uuid{Q8dR}
\exo7id{4435}
\titre{exo7 4435}
\auteur{quercia}
\organisation{exo7}
\datecreate{2010-03-14}
\isIndication{false}
\isCorrection{true}
\chapitre{Série numérique}
\sousChapitre{Autre}
\cpge{mpsi-mp2i.procedes-sommatoires-discrets.a, pcsi.series-numeriques.a, ptsi.series-numeriques.a, mpsi-mp2i.calcul-algebrique-trigonometrie.a, pcsi.calcul-algebrique-trigonometrie.a, ptsi.calcul-algebrique-trigonometrie.a, bcpst1.methodes-calcul, mp-mpi.series-numeriques-vectorielles.a, pt.series-numeriques.b, bcpst2.series-reelles}
\module{Analyse}
\niveau{L1}
\difficulte{}

\contenu{
\question{
Pour quelles valeurs de $a,b\in \R$ la série de terme général
$\ln(n) + a\ln(n+1) + b\ln(n+2)$ est-elle convergente~?
Calculer alors la somme de la série.
}
\reponse{
$\begin{cases} a+b = -1 \cr a+2b = 0\cr \end{cases} \Leftrightarrow a=-2$, $b=1$, $S = -\ln 2$.
}
}
